% !PRINTED 221
\weberpage{221}
%<<CONT
erste Invariante der in Bezug auf die Variablen $\xi_1, \xi_2$ biquadratischen
 Form, n\"amlich (Bd.\ I, \S.\,64):

\weq{3}{u_2^2 - 3\,u_1u_3 + 12\,u_0u_4.}

Da aber eine Form 16\,ten Grades sich nicht als Product von
Formen 12\,ten, 20\,sten und 30\,sten Grades darstellen lassen kann, so
muss diese Covariante identisch verschwinden. Nun ist hier

\[
u_0 = 11\,.\,10\,.\,9\,.\,8\,x_1^7x_2 + 6\,.\,5\,.\,4\,.\,3\,m\,x_1^2x_2^6,
\]
\[
u_1 = 4\,.\,11\,.\,10\,.\,9\,x_1^8 + 4\,.\,6\,.\,5\,.\,4\,.\,6\,m\,x_1^3x_2^5,
\]
\[
u_2 = 6^3\,.\,5^2\,m\,x_1^4x_2^4,
\]
\[
u_3 = -4\,.\,11\,.\,10\,.\,9\,x_2^8 + 4\,.\,6\,.\,5\,.\,4\,.\,6\,m\,x_2^3x_1^5,
\]
\[
u_4 = -11\,.\,10\,.\,9\,.\,8\,x_2^7x_1 + 6\,.\,5\,.\,4\,.\,3\,m\,x_2^2x_1^6,
\]

und wenn wir daraus die Covariante (3) bilden und den Co\"efficienten
 von $x_1^8x_2^8$ gleich 0 setzen, so ergiebt sich $m = \pm 11$.
Beide Zahlen sind hier zul\"assig. Die eine Annahme wird auf
die andere zur\"uckgef\"uhrt durch die Vertauschung von $x_1$ mit
$-x_1$, also durch eine Transformation der Gruppe. Wir setzen
demnach:

\weq{4}{f(x_1,x_2) = x_1x_2(x_1^{10} + 11\,x_1^5x_2^5 - x_2^{10}).}

Die zehn noch fehlenden f\"unfz\"ahligen Pole erh\"alt man also
durch Aufl\"osung der Gleichung

\[
x^{10} + 11\,x^5 - 1 = 0,
\]

aus der man

\[
x^5 = \frac{-11 \pm 5\sqrt5}{2} = \left(\frac{-1 \pm \sqrt5}{2}\right)^5
\]

findet. Setzen wir demnach

\weq{5}{\begin{aligned}
\omega &= \varepsilon + \varepsilon^4 = 2\cos\frac{2\pi}{5} = \frac{-1+\sqrt5}{2} \\
\omega' &= \varepsilon^2 + \varepsilon^3 = 2\cos\frac{4\pi}{5} = \frac{-1-\sqrt5}{2},
\end{aligned}}

so dass

\weq{6}{\omega^2 + \omega = 1, \quad \omega + \omega' = -1, \quad \omega\omega' = -1}

ist, so sind die f\"unfz\"ahligen Pole ausser 0 und $\infty$:

\weq{7}{\xi = \varepsilon^\lambda\omega, \quad \varepsilon^\lambda\omega' = \frac{-\varepsilon^\lambda}{\omega}, \quad \lambda = 0, 1, 2, 3, 4.}

Nun wenden wir den Satz \S.\,53, 2.\ an, nach dem, wenn $a, a'$
zwei Pole einer Substitution $\Theta$ sind, und $b$ ein zu $a$ conjugirter
Pol ist, eine Substitution $\chi$ in der Gruppe existiren muss, so
%CONT>>
